% Study sheet: Volume IV, Chapter 14, Supervisory Intervention.
% Numbers: \val{...} from intervention_studysheet_numbers.tex, written by intervention_studysheet_numbers.py.
\documentclass{studysheet}
\input{intervention_studysheet_numbers}
\sheetvolume{4}
\sheetchapter{14}{Supervisory Intervention}
\sheetsources{Numbers from the chapter's mlsysim scenarios for the warehouse mobile manipulator (illustrative design budgets, not measured takeovers); questions marked assumed use new inputs.}

\begin{document}
% ===================================================================== FRONT: LEARN
\frontside
\begin{bigidea}
A person held in reserve is a delay, not a brake. While the operator responds, the machine
keeps moving, so a takeover is \hl{priced in the stopping budget} with its own deadline, limits and record. Momentum delivered across the \hl{causal boundary} cannot be recalled.
\end{bigidea}

\begin{leftcol}
\sheetsec{Key ideas}
\begin{ideas}
\idea{Intervention moves the right to propose, not to admit.}{A role may propose only on the
  channels its grant names; every command still passes the check.}
\idea{Waiting is travel before braking.}{No brake buys back the wait. Price the human at a
  design percentile, not the mean.}
\idea{The timeout is a function of state.}{Speed and clearance set the latest fallback time; slow down
  before asking, or stop first.}
\idea{Every authority holds a lease.}{Silence expires it into a state-matched fallback: a moving base
  brakes, a loaded arm holds.}
\idea{One writer per channel, announced per channel.}{The enforcer alone writes the actuators;
  show each channel's owner, or hold that channel.}
\idea{The record says who drove each tick.}{Proposed, admitted, dispatched and measured, gaps declared.
  Only stabilized human recovery trains as demonstration.}
\end{ideas}

\vspace{8pt}
\fcard{Bumpless transfer window}%
  {\BumplessTransferBlend = 10\BumplessTransferProgress^3 - 15\BumplessTransferProgress^4 + 6\BumplessTransferProgress^5,
   \qquad \BlendingTransitionWindow \ge \frac{\val{smax}\,|\FiniteChange\ChannelCommand|}{\CommandRateCeiling}}%
  {Progress $\BumplessTransferProgress$ runs from zero to one between latched endpoints; rate and its
   derivative vanish at both ends. $\FiniteChange\ChannelCommand$: command change; $\CommandRateCeiling$: rate ceiling.}%
  {}{}

\vspace{8pt}
\sheetsec{In-loop session at \val{vin}\,m/s: what must fit}
\begin{sheettable}{@{}L r r@{}}
\thead{Charge} & \thead{Travel, mm} & \thead{Left, mm}\\ \midrule
Spare after the finished stop & & \val{tspare}\\
Takeover, \val{tin}\,ms & \val{ttake} & \val{tleft1}\\
Allocation $\DispatchAdmissionInterval$, \val{wdda}\,ms & \val{tda} & \val{tleft2}\\
\val{mlease}\,ms manual lease replaces \val{chunk}\,ms chunk lease & \val{tlease} & \val{tleft3}\\
Braking blend: beyond the resident stop & \val{blendover} & $-$\val{blendshort}\\
\end{sheettable}
\par\vspace{3pt}{\fontsize{8.6}{11}\selectfont\color{muted}Lease floor: heartbeat spacing, jitter, one lost
heartbeat. Ceiling: extra travel $\SpeedAtDetection(\ManualLease - \ChunkLease)$ must fit what is left.\par}
\end{leftcol}%
\begin{rightcol}
\fcard{Takeover delay}%
  {\SupervisoryTakeoverLatency = \OperatorSensoryTime + \OperatorCognitiveTime + \OperatorMotorTime + \ConsoleCommunicationTime}%
  {In the loop: \val{tin}\,ms;
   monitoring, out of the loop: \val{toolms}\,ms.}%
  {}%
  {\begin{tikzpicture}[sheet]
     \pgfmathsetmacro\xa{\val{sv}}
     \pgfmathsetmacro\xb{\xa+\val{ss}}
     \pgfmathsetmacro\xc{\xb+\val{sc}}
     \pgfmathsetmacro\xd{\xc+\val{sm}}
     \pgfmathsetmacro\xe{\xd+\val{sh}}
     \draw[nbox] (0,0) rectangle (\xa,12);
     \draw[accbox] (\xa,0) rectangle (\xb,12);
     \draw[accbox] (\xb,0) rectangle (\xc,12);
     \draw[accbox] (\xc,0) rectangle (\xd,12);
     \draw[nbox] (\xd,0) rectangle (\xe,12);
     \node[lab] at ({(\val{sleft}+\xa)/2},6) {video};
     \node[lab] at ({(\xa+\xb)/2},6) {perceive};
     \node[lab] at ({(\xb+\xc)/2},6) {decide};
     \node[lab] at ({(\xc+\xd)/2},6) {act};
     \node[small, anchor=west] at ({\xe+2},6) {\val{tin}\,ms};
     \draw[dim, dashed] (\val{sleft},-2) -- (\val{sleft},20);
     \node[small, anchor=south west] at ({\val{sleft}-2},21) {waiting allowed at \val{wv}\,m/s: \val{wleft}\,ms};
     \node[note, anchor=north west] at (0,-4) {to scale; last: handshake};
   \end{tikzpicture}}
\fcard{Latest fallback time}%
  {\begin{gathered}\FallbackDeadline = \frac{\ClearDistance - \FixedClearanceAllowance - \ResidentStopDistance(\SpeedAtDetection)}{\SpeedAtDetection} - \PrebrakeDelay\\[2pt]
   \ResidentStopDistance(\SpeedAtDetection) = \frac{\SpeedAtDetection^2}{2\CredibleBrakingDeceleration} + \SpeedAtDetection\frac{\CredibleBrakingDeceleration}{2\JerkLimit}\end{gathered}}%
  {Longest wait before the machine must start its own stop
   ($\FixedClearanceAllowance$: fixed allowances; $\ResidentStopDistance$: jerk-limited stop). Subtract
   the allocation $\DispatchAdmissionInterval$ for the complete budget.}%
  {Timeout $\HandoverTimeout = \max(0, \min(\HandoverTimeoutCeiling, \FallbackDeadline))$.}{}
\fcard{Takeover speed ceiling}%
  {\begin{gathered}\SpeedBound = \CredibleBrakingDeceleration\left[\sqrt{\FullSpeedTravelTime^2 + \frac{2(\ClearDistance - \FixedClearanceAllowance)}{\CredibleBrakingDeceleration}} - \FullSpeedTravelTime\right]\\[2pt]
   \FullSpeedTravelTime = \PrebrakeDelay + \DispatchAdmissionInterval + \SupervisoryTakeoverLatency + \frac{\CredibleBrakingDeceleration}{2\JerkLimit}\end{gathered}}%
  {$\FullSpeedTravelTime$: time at full speed. In-loop: \val{ceilda}\,m/s; the design slows to \val{vin}\,m/s.}%
  {Holds above $\CredibleBrakingDeceleration^2/\JerkLimit$ = \val{vknee}\,m/s.}{}

\vspace{8pt}
\sheetsec{Authority states, per channel}
{\centering\begin{tikzpicture}[sheet]
  \node[nbox, minimum width=40pt, minimum height=13pt] (A) at (24,40) {Autonomous};
  \node[nbox, minimum width=36pt, minimum height=13pt] (P) at (88,40) {Pending};
  \node[hotbox, minimum width=36pt, minimum height=13pt] (B) at (152,40) {Blend};
  \node[accbox, minimum width=36pt, minimum height=13pt] (M) at (216,40) {Manual};
  \node[nbox, minimum width=200pt, minimum height=12pt] (F) at (120,4) {Fallback: validated hold or stop for the current state};
  \draw[arr] (A) -- (P); \node[small, anchor=south] at (56,49) {request};
  \draw[arr] (P) -- (B); \node[small, anchor=south] at (120,49) {commit};
  \draw[arr] (B) -- (M); \node[small, anchor=south] at (184,49) {blend done};
  \foreach \x in {24,88,152,216}{\draw[arr] (\x,33.5) -- (\x,10);}
  \node[small, fill=white, inner sep=1pt] at (120,22) {lapse or timeout};
\end{tikzpicture}\par}
\vspace{2pt}{\fontsize{8.6}{11}\selectfont\color{muted}Fallback exits only by a reset from rest.\par}
\end{rightcol}

% ===================================================================== BACK: PRACTICE
\backside
\begin{answers}
\ans{Finish}{Resident stop \val{ybrake} + \val{yrampt} = \val{yres}\,mm; (\val{wavail} $-$ \val{yres})/\val{yv} = \val{yper}\,ms,
  less \val{wtau}: \res{\val{ydead}\,ms}. The takeover leaves \val{yafter}\,ms; the \val{wdda}\,ms allocation
  then misses by \res{\val{yshort}\,ms}. Partial budget passes, complete fails: slow to \val{vin}\,m/s.}
\ans{1}{$\FullSpeedTravelTime$ = \val{qb}\,s; root of \val{qrad} is \val{qroot}; \val{wa} $\times$ (\val{qroot} $-$ \val{qb}) =
  \res{\val{qceil}\,m/s}. Below the \val{vknee}\,m/s knee the card's stop is longer than the
  true one, so the bound is safe. Out-of-loop help means stop first, then ask (ceiling card).}
\columnbreak
\ans{2}{\val{smax} $\times$ \val{qdtorque} / \val{qrate} = \res{\val{qblend}\,ms}. A linear ramp's rate jumps at both
  ends, a step in torque rate the drivetrain feels (window card).}
\ans{3}{\val{vin} $\times$ \val{qextra}\,ms = \res{\val{qtravel}\,mm} $>$ \val{tleft2}\,mm: misses by \val{qmiss}\,mm. The floor is heartbeat
  spacing plus jitter and a lost heartbeat: fix the link, or request from a crawl or standstill (lease note).}
\ans{4}{Without the proposal, who asked; without the admitted command, whether the enforcer changed it;
  a dispatched torque is not delivered torque: add measured phase current and encoders (record idea).}
\end{answers}

\begin{leftcol}
\sheetsec{Worked example}
\begin{practicebox}\raggedright\boxtitle{Can the base wait for an in-loop takeover at aisle speed?}
The policy requests help at \val{wv}\,m/s, \val{wclear}\,m from the rack end. $\PrebrakeDelay$ = \val{wtau}\,ms,
$\CredibleBrakingDeceleration$ = \val{wa}\,m/s$^2$, $\JerkLimit$ = \val{wj}\,m/s$^3$ (so
$\CredibleBrakingDeceleration/2\JerkLimit$ = \val{wramp}\,ms), $\FixedClearanceAllowance$ = \val{wover}\,mm,
$\DispatchAdmissionInterval$ = \val{wdda}\,ms. Takeover $\SupervisoryTakeoverLatency$ = \val{tin}\,ms.
\par\vspace{5pt}
{\centering
\begin{tikzpicture}[sheet]
  \pgfmathsetmacro\xa{\val{pdel}}
  \pgfmathsetmacro\xb{\xa+\val{pda}}
  \pgfmathsetmacro\xc{\xb+\val{pres}}
  \pgfmathsetmacro\xd{\xc+\val{pover}}
  \pgfmathsetmacro\xt{\xb+\val{ptake}}
  \pgfmathsetmacro\xu{\xt+\val{pres}}
  \pgfmathsetmacro\xv{\xu+\val{pover}}
  \draw[nbox] (0,22) rectangle (\xb,34);
  \draw[accbox] (\xb,22) rectangle (\xc,34);
  \draw[nbox] (\xc,22) rectangle (\xd,34);
  \node[lab] at ({(\xb+\xc)/2},28) {stop};
  \draw[nbox] (0,4) rectangle (\xb,16);
  \draw[hotbox] (\xb,4) rectangle (\xt,16);
  \draw[accbox] (\xt,4) rectangle (\xu,16);
  \draw[nbox] (\xu,4) rectangle (\xv,16);
  \node[lab, text=white] at ({(\xb+\xt)/2},10) {takeover};
  \node[lab] at ({(\xt+\xu)/2},10) {stop};
  \draw[dim, dashed] (\val{pclear},-2) -- (\val{pclear},40);
  \node[small, anchor=south] at (\val{pclear},40) {rack end};
  \node[note, anchor=north west] at (0,-3) {to scale; outlined: delay, $\DispatchAdmissionInterval$ and $\FixedClearanceAllowance$; ends \val{wstopall}\,mm};
\end{tikzpicture}\par}\vspace{8pt}
\begin{steps}
\item Resident stop: $\val{wv}^2/(2\times\val{wa})$ = \val{wbrake}\,mm plus \val{wv} $\times$ \val{wramp} = \val{wrampt}: \res{\val{wres}\,mm}.
\item Room: \val{wclearmm} $-$ \val{wover} $-$ \val{wres} = \val{wroom}\,mm; $\div$ \val{wv}\,m/s = \val{wper}\,ms; less \val{wtau}: \res{$\FallbackDeadline$ = \val{wdead}\,ms}.
\item Charge $\DispatchAdmissionInterval$: \val{wdead} $-$ \val{wdda} = \res{\val{wleft}\,ms} of waiting.
\item The takeover needs \val{tin}\,ms, about \val{wratio} times that.
\end{steps}
\takeaway{It cannot wait. Stop locally, then ask; or slow before an anticipated request.}
\end{practicebox}

\vspace{10pt}
\sheetsec{Your turn: finish it}
\begin{practicebox}\boxtitle{Same request, after a slow-down to \val{yv}\,m/s}
A designer picks \val{yv}\,m/s, under the \val{ceilpartial}\,m/s ceiling found without $\DispatchAdmissionInterval$.
Everything else is unchanged.\par\vspace{4pt}
\textbf{Given:} resident stop \val{ybrake} + \val{yrampt}\,mm.\par
\textbf{You:} find $\FallbackDeadline$, take off the takeover and $\DispatchAdmissionInterval$, and decide.\par\vspace{30pt}
\end{practicebox}
\end{leftcol}%
\begin{rightcol}
\sheetsec{Common mistakes}
\mistake{``A fixed \val{window}\,ms handover timeout is generous; it is safe everywhere.''}%
  {At \val{wv}\,m/s that wait is \val{fwtravel}\,m of unbraked travel, and the stop ends \val{fwovercomplete}\,m past the
   rack end. At a \val{fwcrawl}\,m/s crawl the fallback deadline is \val{fwmargin}\,s. Compute the deadline
   every tick.}{}
\mistake{``Only the enforcer writes the actuators, so nobody can be confused about who drives.''}%
  {Single writer fixes arbitration, not belief. An operator who spends \val{tin}\,ms aligning the arm while the
   unannounced base runs at \val{wv}\,m/s lets it travel \val{m2travel}\,mm, \val{m2ratio} times the \val{m2spare}\,mm left after dispatch travel.}{}

\vspace{2pt}
\sheetsec{Check yourself}
\begin{qs}
\item \assumed{} the request reaches a monitoring operator, out of the loop: $\SupervisoryTakeoverLatency$ = \val{toolms}\,ms. With the
  worked example's givens, what is the complete-budget ceiling? Compare the in-loop \val{ceilda}\,m/s.
\item \assumed{} the operator's first arm target is \val{qdtorque}\,N$\cdot$m from the latched command;
  torque-rate ceiling \val{qrate}\,N$\cdot$m/s. Minimum blend window? Why not a linear ramp?
\item \assumed{} a lossy link: the site wants a \val{qlease}\,ms manual lease at \val{vin}\,m/s. Does it fit
  the session budget? What sets the lease's lower bound?
\item A log keeps only the authority state and the dispatched torque per tick. After the arm strikes
  a rack upright under manual control, which questions can it not answer, and what would?
\end{qs}
\end{rightcol}
\end{document}
